\documentclass[11pt]{article}

\usepackage[a4paper,margin=1in]{geometry}
\usepackage{amsmath}
\usepackage{amssymb}
\usepackage{booktabs}
\usepackage{hyperref}
\usepackage{xcolor}
\usepackage{listings}

\lstset{
  basicstyle=\ttfamily\small,
  breaklines=true,
  frame=single,
  columns=fullflexible
}

\title{Robot-and-Object GT-Dynamics-Supervised Hybrid DeLaN World Model}
\author{}
\date{}

\begin{document}
\maketitle

\section{Overview}

The implementation under
\begin{center}
\texttt{scripts/world\_model/Physics/Robot\_and\_Object/GT\_dynamics}
\end{center}
trains a torque-driven hybrid world model for Franka lift episodes. It extends
the robot-only GT-dynamics model by including the manipulated object's state,
object physical context, object dynamics labels, and a learned residual contact
coupling term.

The robot part remains DeLaN-style and predicts structured robot dynamics terms:
\begin{equation}
  H(q), \qquad c(q,\dot{q}), \qquad g(q).
\end{equation}
The object/contact part is learned with MLPs conditioned on the full robot-object
state, recorded robot torque, and object context. It predicts
\begin{align}
  \tau_{\mathrm{res}} &= f_{\theta}^{\tau}(s_t,\tau_t,\xi), \\
  a_o &= f_{\theta}^{o}(s_t,\tau_t,\xi),
\end{align}
where \(\tau_{\mathrm{res}}\) is a residual generalized torque and \(a_o\)
contains object linear and angular accelerations.

This design is hybrid rather than a single full-system DeLaN model. The HDF5
dataset contains clean robot dynamics labels and rigid-object state/force labels,
but it does not contain a clean full generalized robot-object mass matrix with
contact constraints. Therefore, the implementation preserves the structured
robot dynamics and learns contact/object effects as additional supervised terms.

\section{State, Input, and Context}

The model uses \(n=9\) robot degrees of freedom by default: seven Franka arm
joints and two gripper finger joints. The robot state is
\begin{equation}
  s^r_t =
  \begin{bmatrix}
    q_t \\
    \dot{q}_t
  \end{bmatrix}
  \in \mathbb{R}^{18}.
\end{equation}

The object state is
\begin{equation}
  s^o_t =
  \begin{bmatrix}
    p^o_t \\
    r^o_t \\
    v^o_t \\
    \omega^o_t
  \end{bmatrix}
  \in \mathbb{R}^{13},
\end{equation}
where \(p^o_t \in \mathbb{R}^3\) is object position,
\(r^o_t \in \mathbb{R}^4\) is a unit quaternion in \texttt{wxyz} order,
\(v^o_t \in \mathbb{R}^3\) is linear velocity, and
\(\omega^o_t \in \mathbb{R}^3\) is angular velocity.

The full state is
\begin{equation}
  s_t =
  \begin{bmatrix}
    q_t \\
    \dot{q}_t \\
    p^o_t \\
    r^o_t \\
    v^o_t \\
    \omega^o_t
  \end{bmatrix}
  \in \mathbb{R}^{31}.
\end{equation}

The dynamics input is the recorded robot torque
\begin{equation}
  \tau_t \in \mathbb{R}^{9}.
\end{equation}
The torque is not predicted by the model. It is read from the dataset using
\texttt{robot\_torques/applied\_torque} or
\texttt{robot\_torques/computed\_torque}.

The model also receives an object context vector
\begin{equation}
  \xi =
  \begin{bmatrix}
    m_o \\
    I_o \\
    \mu_o
  \end{bmatrix}
  \in \mathbb{R}^{13},
\end{equation}
where \(m_o \in \mathbb{R}\) is object mass,
\(I_o \in \mathbb{R}^9\) is the recorded inertia vector, and
\(\mu_o \in \mathbb{R}^3\) is the recorded material-property vector.

\section{Expected Dataset Fields}

Each HDF5 episode is expected to contain robot observations and recorded torque:
\begin{lstlisting}
obs/joint_pos
obs/joint_vel
actions
robot_torques/applied_torque
robot_torques/computed_torque
\end{lstlisting}

The robot GT dynamics labels are:
\begin{lstlisting}
robot_dynamics/qdd
robot_dynamics/mass_matrix
robot_dynamics/inertial
robot_dynamics/coriolis
robot_dynamics/gravity
robot_dynamics/inverse_dynamics_tau
\end{lstlisting}

The object state and object dynamics labels are:
\begin{lstlisting}
states/rigid_object/object/root_pose
states/rigid_object/object/root_velocity
object_dynamics/root_lin_acc_w
object_dynamics/root_ang_acc_w
object_dynamics/mass
object_dynamics/inertia
object_dynamics/material_properties
object_dynamics/inertial_force_w
object_dynamics/gravity_force_w
object_dynamics/external_force_est_w
\end{lstlisting}

The current default dataset path in \texttt{wm\_train.py} is:
\begin{lstlisting}
../../../../reinforcement_learning/skrl/datasets/Lift_RL_opt_robot_object_dynamics_10000ep.hdf5
\end{lstlisting}

\section{Training Samples}

The dataset class is
\begin{center}
\texttt{RobotObjectGTDynamicsRolloutDataset}.
\end{center}
It builds sliding-window samples without crossing episode boundaries. Each sample
contains
\begin{align}
  \texttt{history\_states} &\in \mathbb{R}^{K_h \times 31}, \\
  \texttt{future\_torques} &\in \mathbb{R}^{H \times 9}, \\
  \texttt{future\_states} &\in \mathbb{R}^{H \times 31}, \\
  \texttt{object\_context} &\in \mathbb{R}^{13}.
\end{align}
Here \(K_h\) is \texttt{history\_len} and \(H\) is
\texttt{rollout\_horizon}.

The loader also returns:
\begin{lstlisting}
future_robot_dynamics
future_object_dynamics
history_actions
future_actions
\end{lstlisting}
The action tensors are kept for compatibility and inspection, but the current
model uses recorded torque as its dynamics input.

Unlike the robot-only pre-contact experiment, the default robot-and-object
trainer keeps contact windows:
\begin{lstlisting}
filter_pre_contact = False
\end{lstlisting}
This is intentional because the goal of this model is to reduce the large error
that appears when the robot contacts and carries the object. A pre-contact-only
filter can still be enabled with
\begin{lstlisting}[language=bash]
python wm_train.py --filter_pre_contact
\end{lstlisting}

\section{Model Input and Output}

The main rollout model is
\begin{center}
\texttt{MultiStepRobotObjectGTDynamicsWorldModel}.
\end{center}
Its forward input is
\begin{align}
  \texttt{history\_states} &\in \mathbb{R}^{B \times K_h \times 31}, \\
  \texttt{future\_torques} &\in \mathbb{R}^{B \times H \times 9}, \\
  \texttt{object\_context} &\in \mathbb{R}^{B \times 13}.
\end{align}

As in the robot-only GT-dynamics model, the current implementation uses the most
recent state in the history:
\begin{equation}
  \hat{s}_t = \texttt{history\_states}_{:, -1}.
\end{equation}
The future torques are taken from the dataset and supplied step by step.

The output is an open-loop predicted rollout:
\begin{equation}
  \hat{S}_{t+1:t+H}
  =
  \left[
    \hat{s}_{t+1},
    \hat{s}_{t+2},
    \ldots,
    \hat{s}_{t+H}
  \right]
  \in \mathbb{R}^{B \times H \times 31}.
\end{equation}

If \texttt{return\_aux=True}, the model also returns auxiliary predictions:
\begin{lstlisting}
H
L
l_diag
l_offdiag
g
coriolis
torque
tau_residual
effective_torque
ddq
inertial
inverse_dynamics_tau
object_lin_acc
object_ang_acc
object_inertial_force
object_gravity_force
object_external_force
\end{lstlisting}

\section{Robot DeLaN-Style Structured Dynamics}

The robot structured terms are computed by
\begin{center}
\texttt{DeLaNStructuredTerms}.
\end{center}
The implementation is the same basic construction as the robot-only GT-dynamics
model. The network directly outputs
\begin{align}
  g(q) &\in \mathbb{R}^{n}, \\
  l_d(q) &\in \mathbb{R}^{n}, \\
  l_o(q) &\in \mathbb{R}^{n(n-1)/2}.
\end{align}
The diagonal output is passed through softplus:
\begin{equation}
  L_{ii}(q) = \mathrm{softplus}(l_{d,i}(q)) + \epsilon_d.
\end{equation}
The inertia matrix is constructed as
\begin{equation}
  H(q) = L(q)L(q)^\top + \epsilon I,
\end{equation}
which enforces positive definiteness up to numerical precision.

The Coriolis and centrifugal term is computed from the learned inertia structure:
\begin{equation}
  c(q,\dot{q})
  =
  \dot{H}(q,\dot{q})\dot{q}
  -
  \frac{1}{2}
  \nabla_q
  \left(
    \dot{q}^{\top}H(q)\dot{q}
  \right).
\end{equation}
The implementation uses a custom \texttt{DerivativeMLP} to propagate
\(\partial L/\partial q\) through the network.

\section{Residual Contact Torque}

The robot-only model solves
\begin{equation}
  H(q_t)\ddot{q}_t =
  \tau_t - c(q_t,\dot{q}_t) - g(q_t).
\end{equation}
This ignores contact and object-carrying effects except through the recorded
robot motion. The robot-and-object model introduces a learned residual
generalized torque:
\begin{equation}
  \tau_{\mathrm{eff},t}
  =
  \tau_t
  +
  \tau_{\mathrm{res},t}.
\end{equation}
The residual is predicted from the full state, torque, and object context:
\begin{equation}
  \tau_{\mathrm{res},t}
  =
  f_{\theta}^{\tau}
  \left(
    s_t,
    \tau_t,
    \xi
  \right).
\end{equation}

The robot forward dynamics equation becomes
\begin{equation}
  H(q_t)\ddot{q}_t
  =
  \tau_t
  +
  \tau_{\mathrm{res},t}
  -
  c(q_t,\dot{q}_t)
  -
  g(q_t).
\end{equation}
In code, this is implemented as:
\begin{lstlisting}[language=Python]
effective_torque = torque + tau_residual
rhs = (effective_torque - g - coriolis).unsqueeze(-1)
ddq = torch.linalg.solve(H, rhs).squeeze(-1)
\end{lstlisting}

The residual target used by the supervised loss is
\begin{equation}
  \tau_{\mathrm{res,GT}}
  =
  \tau_{\mathrm{ID,GT}}
  -
  \tau_t,
\end{equation}
where
\begin{equation}
  \tau_{\mathrm{ID,GT}}
  =
  M_{\mathrm{GT}}\ddot{q}_{\mathrm{GT}}
  +
  c_{\mathrm{GT}}
  +
  g_{\mathrm{GT}}.
\end{equation}
This target should be interpreted as a useful residual supervision signal, not
as a perfect measured contact force. It includes whatever mismatch exists
between recorded actuator torque and the PhysX inverse-dynamics label.

\section{Object Dynamics and Integration}

The object head predicts six values:
\begin{equation}
  a^o_t =
  \begin{bmatrix}
    \dot{v}^o_t \\
    \dot{\omega}^o_t
  \end{bmatrix}
  =
  f_{\theta}^{o}(s_t,\tau_t,\xi)
  \in \mathbb{R}^{6}.
\end{equation}
The first three dimensions are object linear acceleration and the last three are
object angular acceleration.

The object velocity and position are integrated with semi-implicit Euler:
\begin{align}
  v^o_{t+1} &= v^o_t + \dot{v}^o_t \Delta t, \\
  p^o_{t+1} &= p^o_t + v^o_{t+1}\Delta t, \\
  \omega^o_{t+1} &= \omega^o_t + \dot{\omega}^o_t \Delta t.
\end{align}

The quaternion is updated using the predicted angular velocity:
\begin{equation}
  \dot{r}^o_t
  =
  \frac{1}{2}
  \begin{bmatrix}
    0 \\
    \omega^o_{t+1}
  \end{bmatrix}
  \otimes
  r^o_t,
\end{equation}
followed by Euler integration and normalization:
\begin{equation}
  r^o_{t+1}
  =
  \mathrm{normalize}
  \left(
    r^o_t + \dot{r}^o_t \Delta t
  \right).
\end{equation}

The model also computes predicted object forces:
\begin{align}
  \hat{F}_{\mathrm{inertial}} &= m_o \hat{\dot{v}}^o, \\
  \hat{F}_{g} &= m_o
  \begin{bmatrix}
    0 & 0 & -9.81
  \end{bmatrix}^{\top}, \\
  \hat{F}_{\mathrm{ext}} &=
  \hat{F}_{\mathrm{inertial}} - \hat{F}_{g}.
\end{align}
The external force prediction is supervised against
\texttt{object\_dynamics/external\_force\_est\_w}.

\section{Open-Loop Multi-Step Rollout}

The model is trained recursively in open loop. The initial state is the real
dataset state:
\begin{equation}
  \hat{s}_t = s_t.
\end{equation}
For each future step \(h=0,\ldots,H-1\),
\begin{equation}
  \hat{s}_{t+h+1}
  =
  f_{\theta}
  \left(
    \hat{s}_{t+h},
    \tau_{t+h},
    \xi
  \right).
\end{equation}
The future torques are always taken from the HDF5 dataset. After the first step,
the model feeds its own predicted robot and object state back into the next
dynamics step.

\section{Training Objective}

The total training loss is
\begin{equation}
  \mathcal{L}
  =
  \mathcal{L}_{\mathrm{rollout}}
  +
  \mathcal{L}_{\mathrm{robotGT}}
  +
  \mathcal{L}_{\mathrm{objectGT}}
  +
  \mathcal{L}_{\mathrm{reg}}
  +
  \mathcal{L}_{\mathrm{L2}}.
\end{equation}

\subsection{Rollout State Loss}

The rollout loss compares predicted future states against ground truth:
\begin{align}
  \mathcal{L}_{\mathrm{rollout}}
  &=
  w_q
  \left\|
    \hat{q} - q
  \right\|_2^2
  +
  w_{\dot{q}}
  \left\|
    \hat{\dot{q}} - \dot{q}
  \right\|_2^2
  +
  w_p
  \left\|
    \hat{p}^o - p^o
  \right\|_2^2 \nonumber \\
  &\quad+
  w_r
  d_r(\hat{r}^o,r^o)
  +
  w_v
  \left\|
    \hat{v}^o - v^o
  \right\|_2^2
  +
  w_{\omega}
  \left\|
    \hat{\omega}^o - \omega^o
  \right\|_2^2.
\end{align}

The quaternion loss uses the sign-invariant distance:
\begin{equation}
  d_r(\hat{r},r)
  =
  \min
  \left(
    \left\|\hat{r} - r\right\|_2^2,
    \left\|\hat{r} + r\right\|_2^2
  \right),
\end{equation}
because \(r\) and \(-r\) represent the same orientation.

The default rollout weights are:
\begin{align}
  w_q &= 1.0, &
  w_{\dot{q}} &= 0.2, &
  w_p &= 5.0, \\
  w_r &= 0.5, &
  w_v &= 1.0, &
  w_{\omega} &= 0.5.
\end{align}

\subsection{Robot GT Dynamics Loss}

The robot dynamics loss supervises the same structured robot terms as the
robot-only version, plus the residual torque:
\begin{align}
  \mathcal{L}_{\mathrm{robotGT}}
  &=
  \lambda_M
  \left\|
    \hat{H} - M_{\mathrm{GT}}
  \right\|_2^2
  +
  \lambda_I
  \left\|
    \widehat{H\ddot{q}} - (M\ddot{q})_{\mathrm{GT}}
  \right\|_2^2 \nonumber \\
  &\quad+
  \lambda_C
  \left\|
    \hat{c} - c_{\mathrm{GT}}
  \right\|_2^2
  +
  \lambda_g
  \left\|
    \hat{g} - g_{\mathrm{GT}}
  \right\|_2^2 \nonumber \\
  &\quad+
  \lambda_{\ddot{q}}
  \left\|
    \hat{\ddot{q}} - \ddot{q}_{\mathrm{GT}}
  \right\|_2^2
  +
  \lambda_{\mathrm{ID}}
  \left\|
    \hat{\tau}_{\mathrm{ID}} - \tau_{\mathrm{ID,GT}}
  \right\|_2^2 \nonumber \\
  &\quad+
  \lambda_{\mathrm{res}}
  \left\|
    \hat{\tau}_{\mathrm{res}} -
    \left(
      \tau_{\mathrm{ID,GT}} - \tau
    \right)
  \right\|_2^2.
\end{align}

The default robot dynamics weights are:
\begin{align}
  \lambda_M &= 0.01, &
  \lambda_I &= 0.01, &
  \lambda_C &= 0.01, \\
  \lambda_g &= 0.01, &
  \lambda_{\ddot{q}} &= 0.0, &
  \lambda_{\mathrm{ID}} &= 0.0, \\
  \lambda_{\mathrm{res}} &= 0.01.
\end{align}

\subsection{Object Dynamics Loss}

The object dynamics loss supervises predicted linear acceleration, angular
acceleration, and external force:
\begin{align}
  \mathcal{L}_{\mathrm{objectGT}}
  &=
  \lambda_{\dot{v}}
  \left\|
    \hat{\dot{v}}^o - \dot{v}^o_{\mathrm{GT}}
  \right\|_2^2
  +
  \lambda_{\dot{\omega}}
  \left\|
    \hat{\dot{\omega}}^o - \dot{\omega}^o_{\mathrm{GT}}
  \right\|_2^2 \nonumber \\
  &\quad+
  \lambda_F
  \left\|
    \hat{F}_{\mathrm{ext}} -
    F_{\mathrm{ext,GT}}
  \right\|_2^2.
\end{align}

The default object dynamics weights are:
\begin{align}
  \lambda_{\dot{v}} &= 0.01, \\
  \lambda_{\dot{\omega}} &= 0.001, \\
  \lambda_F &= 0.01.
\end{align}

\subsection{Regularization}

Two auxiliary regularizers are available:
\begin{equation}
  \mathcal{L}_{\mathrm{reg}}
  =
  \lambda_{\mathrm{cor}}
  \left\|
    \hat{c}
  \right\|_2^2
  +
  \lambda_{\tau}
  \left\|
    \hat{\tau}_{\mathrm{res}}
  \right\|_2^2.
\end{equation}
The default values are:
\begin{align}
  \lambda_{\mathrm{cor}} &= 0.0, \\
  \lambda_{\tau} &= 10^{-5}.
\end{align}

The trainer also applies an explicit L2 penalty to trainable weight matrices:
\begin{equation}
  \mathcal{L}_{\mathrm{L2}}
  =
  \lambda_{\mathrm{W}}
  \sum_i
  \|W_i\|_2^2,
\end{equation}
with default
\begin{equation}
  \lambda_{\mathrm{W}} = 10^{-6}.
\end{equation}

\section{Training Procedure}

Training uses AdamW:
\begin{align}
  \text{learning rate} &= 10^{-5}, \\
  \text{weight decay} &= 10^{-6}, \\
  \text{batch size} &= 256.
\end{align}
The default rollout horizon is
\begin{equation}
  H = 3,
\end{equation}
and the default history length is
\begin{equation}
  K_h = 5.
\end{equation}

The dataset is split by episode:
\begin{equation}
  90\% \text{ training episodes}, \qquad
  10\% \text{ validation episodes}.
\end{equation}
For each epoch, the model is trained on mini-batches of rollout windows. The
validation loss is computed without optimizer updates. The trainer saves:
\begin{lstlisting}
best.pt
last.pt
\end{lstlisting}
where \texttt{best.pt} is selected using the lowest validation total loss.

Gradient clipping is applied:
\begin{equation}
  \|\nabla_\theta \mathcal{L}\|_2 \leq 10.
\end{equation}

\section{Logged Metrics}

The training loop logs the following major quantities to TensorBoard:
\begin{lstlisting}
loss
rollout_loss
robot_dyn_loss
object_dyn_loss
reg_loss
weight_l2_loss
q_mse
dq_mse
object_pos_mse
object_quat_mse
object_lin_vel_mse
object_ang_vel_mse
one_step_q_mse
one_step_object_pos_mse
final_step_q_mse
final_step_object_pos_mse
mass_matrix_mse
inertial_mse
coriolis_gt_mse
gravity_mse
qdd_mse
inverse_dynamics_mse
residual_torque_mse
object_lin_acc_mse
object_ang_acc_mse
object_external_force_mse
coriolis_reg
residual_torque_reg
torque_abs_mean
residual_torque_abs_mean
\end{lstlisting}

\section{Example Commands}

Training with the default hybrid loss:
\begin{lstlisting}[language=bash]
python wm_train.py \
  --dataset_file ../../../../reinforcement_learning/skrl/datasets/Lift_RL_opt_robot_object_dynamics_10000ep.hdf5 \
  --rollout_horizon 3 \
  --train_batch_fraction 0.1
\end{lstlisting}

Training with stronger object acceleration supervision:
\begin{lstlisting}[language=bash]
python wm_train.py \
  --dataset_file ../../../../reinforcement_learning/skrl/datasets/Lift_RL_opt_robot_object_dynamics_10000ep.hdf5 \
  --lambda_object_lin_acc 0.05 \
  --lambda_object_ang_acc 0.005 \
  --lambda_object_external_force 0.05
\end{lstlisting}

Plotting robot gripper and object trajectories:
\begin{lstlisting}[language=bash]
python plot_robot_object_trajectory.py \
  --checkpoint ./outputs_robot_object_gt_dynamics/run_YYYYMMDD_HHMMSS/best.pt \
  --dataset_file ../../../../reinforcement_learning/skrl/datasets/Lift_RL_opt_robot_object_dynamics_10000ep.hdf5 \
  --episode_index 0 \
  --rollout_steps 100
\end{lstlisting}

Plotting robot/object dynamics terms:
\begin{lstlisting}[language=bash]
python plot_gt_vs_predicted_dynamics_terms.py \
  --checkpoint ./outputs_robot_object_gt_dynamics/run_YYYYMMDD_HHMMSS/best.pt \
  --dataset_file ../../../../reinforcement_learning/skrl/datasets/Lift_RL_opt_robot_object_dynamics_10000ep.hdf5 \
  --episode_index 0 \
  --start_t 20 \
  --rollout_steps 50
\end{lstlisting}

\section{Important Caveats}

\subsection{Torque Is Still an Input}

The model receives recorded future torque values from the dataset:
\begin{equation}
  \hat{\tau}_t = \tau_t.
\end{equation}
It does not predict policy actions or actuator torques. The learned residual
torque is an additional correction used inside the dynamics equation, not a
replacement for the recorded torque input.

\subsection{Robot Dynamics Labels Are Free-Space Terms}

The labels \(M\), \(c\), and \(g\) come from PhysX robot dynamics queries. They
are free-space robot model terms. Contact and grasp effects are represented in
this architecture through the learned residual torque and object dynamics
supervision, not through a full constrained robot-object mass matrix.

\subsection{Residual Torque Is Not a Direct Force Sensor}

The residual target
\begin{equation}
  \tau_{\mathrm{ID,GT}} - \tau
\end{equation}
is useful for learning mismatch around contact, but it is not a direct measured
contact wrench. It can include actuator-model mismatch, numerical effects, and
differences between recorded torque and inverse-dynamics labels.

\subsection{Object Force Estimate Is Approximate}

The dataset's \texttt{external\_force\_est\_w} is computed from object mass,
linear acceleration, and gravity:
\begin{equation}
  F_{\mathrm{ext}} \approx m_o(\dot{v}^o - g).
\end{equation}
It is an estimated non-gravity force in world coordinates. It can reflect
contact and grasp effects, but it is not a direct contact-force measurement.

\subsection{Open-Loop Instability}

During long open-loop rollout, the model recursively feeds predicted robot and
object states back into the next dynamics step. Errors can accumulate:
\begin{equation}
  \hat{s}_{t+h}
  \neq
  s_{t+h}.
\end{equation}
If the predicted robot or object state leaves the training distribution,
structured terms, residual torques, object accelerations, or quaternion
integration can become inaccurate. The evaluation scripts stop and report the
failure if non-finite states appear.

\subsection{No Full Coupled DeLaN System}

This model is not a full DeLaN model over robot and object generalized
coordinates. It does not construct a coupled mass matrix for
\([q, p^o, r^o]\). Instead, it keeps the robot dynamics structured and learns
object/contact coupling with supervised residual heads. This is a practical
choice matched to the available HDF5 labels.

\end{document}
